\section{System Model}
\label{sec:system}

% figures/tikz_architecture.tex  ->  \cref{fig:architecture}
\begin{figure}[t]
\centering
\resizebox{\textwidth}{!}{%
\begin{tikzpicture}[node distance=8mm and 13mm]

% ---- lane 1: honest data path (left to right) ----
\node[payload] (payload) {\icondoc[aegisAmber]\\ untrusted\\ payload $\pay$};
\node[honest, right=of payload] (input) {\iconserver[aegisBlue]\\ Input\\ agent};
\node[honest, right=of input] (analyzer)
  {\iconmodel[aegisBlue]\\ Analyzer\\ agents};

% ---- isolation gate (defense) ----
\node[defense, right=16mm of analyzer] (iso)
  {\iconlock[aegisGreen]\\ payload\\ isolation};

% ---- judge committee (stacked); Byzantine judge placed at top so the
%      compromise arrow drops vertically without crossing other judges ----
\node[byz,    right=18mm of iso, yshift=15mm] (j3) {\iconjudge[aegisRed]\ Judge $k$ (Byz.)};
\node[honest, below=5mm of j3] (j1) {\iconjudge[aegisBlue]\ Judge $1$};
\node[honest, below=5mm of j1] (j2) {\iconjudge[aegisBlue]\ Judge $2$};
\node[honest, below=5mm of j2] (jn) {\iconjudge[aegisBlue]\ Judge $n$};

% committee grouping box
\begin{scope}[on background layer]
  \node[group, fit=(j3)(j1)(j2)(jn),
        label={[grouplabel, xshift=16mm]above:hardened committee ($n$)}] (comm) {};
\end{scope}

% ---- robust aggregator (defense) ----
\node[defense, right=18mm of comm] (agg)
  {\iconshield[aegisGreen]\\ Byzantine-robust\\ aggregation\\ $\Agg\in\{\Krum,\cmed,\gmed\}$};
\node[honest, right=of agg] (out)
  {\iconserver[aegisBlue]\\ Output\\ agent};
\node[theory, below=9mm of out] (dec)
  {decision\\ $\hat d=\dec(\uhat)$};

% ---- attacker (adversarial lane, top), placed above the Byzantine judge ----
\node[adv, above=9mm of j3] (atk) {\iconattacker[aegisRed]\\ adaptive\\ attacker};

% ===== flows =====
% data/model flow (blue)
\draw[dataflow] (payload) -- (input);
\draw[dataflow] (input) -- (analyzer);
\draw[dataflow] (analyzer) -- (iso);
\draw[dataflow] (iso.east) -- (j1.west);
\draw[dataflow] (iso.east) -- (j2.west);
\draw[dataflow] (iso.east) -- (j3.west);
\draw[dataflow] (iso.east) -- (jn.west);
\draw[dataflow] (j1.east) -- (agg.west);
\draw[dataflow] (j2.east) -- (agg.west);
\draw[dataflow] (j3.east) -- (agg.west);
\draw[dataflow] (jn.east) -- (agg.west);
\draw[dataflow] (agg) -- (out);
\draw[dataflow] (out) -- (dec);

% adversarial flow (red): inject payload (routed along the top) + compromise judge k
\draw[advflow] (atk.north) |- ($(payload.north)+(0,10mm)$) -| (payload.north)
  node[elabel,pos=0.25]{second-order injection};
\draw[advflow] (atk.south) -- (j3.north)
  node[elabel,pos=0.55,xshift=-11mm]{compromise};

% control/verdict labels on chips
\node[elabel] at ($(j1.east)!0.5!(agg.west)+(0,3mm)$) {verdicts $\verd_k$};
\node[elabel] at ($(iso.east)!0.5!(j2.west)$) {data ch.};

\end{tikzpicture}}
\caption{The \Aegis\ architecture. Blue: honest data/model flow through Input, Analyzer,
isolation, the hardened judge committee, robust aggregation, and Output. Green:
defensive components (payload isolation and Byzantine-robust verdict aggregation) that
replace AutoDefense's single Coordinator. Red (dashed): the adaptive attacker, which both
injects the untrusted payload (second-order injection against the judges) and compromises
up to $f$ judges (here Judge~$k$). Judges do not communicate with one another; each sends
only its verdict $\verd_k$ to the aggregator. Icons are pure TikZ. Arrow lanes are separated
by colour and dash pattern; labels sit on white chips so no connector crosses text.}
\label{fig:architecture}
\end{figure}


\subsection{The defense agency}
\Aegis is a defensive multi-agent system that adjudicates whether a candidate output may
be released. It comprises (\cref{fig:architecture}):
\begin{itemize}[leftmargin=1.4em,itemsep=2pt,topsep=2pt]
  \item an \textbf{Input agent} that receives the candidate output and any accompanying
        context and places the untrusted content into a structured data channel;
  \item one or more \textbf{Analyzer agents} that extract task-relevant features—the
        apparent user intent and the original instruction/privilege
        hierarchy~\cite{wallace2024instruction}—and attach them to the instruction
        channel as trusted metadata;
  \item a committee of $n$ \textbf{Judge agents} ($1\le n\le 7$ in our study), each a
        hardened LLM~\cite{chen2025secalign,chen2024struq} that reads the isolated
        payload and emits a verdict $v_k$ as in \cref{eq:verdict};
  \item a \textbf{robust aggregator} that replaces AutoDefense's single Coordinator with
        a Byzantine-robust rule $\Agg$ over $\{\verd_k\}$; and
  \item an \textbf{Output agent} that enforces the aggregate decision $\hat d$ (release,
        block, or escalate).
\end{itemize}
The committee is deployed in one of two topologies. In the \emph{parallel} topology the
$n$ judges are queried independently and concurrently; in the \emph{sequential} topology
they are queried in series (relevant only to latency, \cref{sec:complexity}). Judges do
not exchange messages with one another; each communicates only with the aggregator. This
star topology is deliberate: it removes the inter-agent communication channel that
recent attacks exploit~\cite{he2025redteaming,lee2024promptinfection} and confines
adjudication to verdicts the aggregator can robustly combine.

\subsection{Trust boundaries}
Three assets are \emph{trusted}: the aggregation substrate executes $\Agg$ correctly and
non-interactively; the instruction channel delivering the operator's policy is authentic;
and at least $n-f$ judges are honest, with $f$ small enough to satisfy the tolerance
condition of the chosen rule ($2f{+}2<n$ for $\Krum$, $f<n/2$ for median rules). Two
assets are \emph{untrusted}: the payload $\pay$, which may be adversarially crafted, and
up to $f$ judges, which may be compromised or colluding. The aggregator being trusted is
what distinguishes our setting from Byzantine agreement among mutually distrustful
parties~\cite{lamport1982byzantine,castro1999pbft}: we need robust \emph{estimation} of a
verdict, not a consensus protocol. This mirrors the trusted-server assumption of robust
federated aggregation~\cite{blanchard2017krum,yin2018byzantine}, transplanted from
gradients to verdicts.

\subsection{Threat model}
\label{sec:threatmodel}
We adopt the standard security-analysis structure of goal, knowledge, and capability.

\paragraph{Attacker goal.}
Two objectives, dual to each other:
\begin{itemize}[leftmargin=1.4em,itemsep=1pt,topsep=2pt]
  \item \emph{Integrity violation:} cause the agency to \textbf{pass} an output that is
        unsafe or that violates the intended task/instruction
        hierarchy—an evasion, measured by attack success rate under compromise.
  \item \emph{Availability violation:} cause the agency to \textbf{block} benign
        outputs—over-refusal, a denial-of-service against legitimate
        traffic~\cite{zhou2025corba}.
\end{itemize}

\paragraph{Attacker knowledge.}
Black-box to gray-box. The attacker may know the pipeline design, the aggregation rule,
the judges' backbones, and the policy prompt, but not the operator's private randomness
(e.g.\ any per-request nonce used for isolation delimiting). This spans the practical
range from a third party who only submits payloads to an insider who fine-tuned a judge.

\paragraph{Attacker capabilities.}
Three levers, usable in combination:
\begin{enumerate}[leftmargin=1.6em,itemsep=1pt,topsep=2pt]
  \item \textbf{Insider compromise.} Control up to $f$ judges, which then return
        arbitrary, jointly optimised verdict vectors. This is the LLM-committee analogue
        of Byzantine workers~\cite{blanchard2017krum} and of poisoned federated
        clients~\cite{bagdasaryan2020backdoor,fang2020poisoning}.
  \item \textbf{Second-order injection.} Craft the payload $\pay$ so that, in addition to
        being the object under review, it carries instructions aimed at the honest judges
        reading it~\cite{shi2024judgeinjection,greshake2023injection}. Without isolation
        this can flip honest verdicts; with isolation it is reduced to the leakage
        $\varepsilon$ of \cref{def:isolation}.
  \item \textbf{Adaptivity.} Optimise the attack against the specific aggregation rule,
        as adaptive attacks do against robust federated
        aggregators~\cite{fang2020poisoning,baruch2019little,xie2019fall}. Colluding
        judges may coordinate to stay within the honest verdict spread while biasing the
        aggregate.
\end{enumerate}

\paragraph{Collusion.}
Byzantine judges are not independent: they may share a strategy and coordinate their
verdict vectors to move the aggregate as far as the tolerance condition permits. Our
analysis and evaluation treat collusion as the default, not the worst case of an
otherwise-independent fault.

\paragraph{Out of scope.}
We do not defend against (i)~all judges compromised ($f\ge n-f$), which no
minority-tolerance rule can address; (ii)~a correlated failure that affects \emph{every}
judge identically—studied in \cref{sec:theory} as a limiting case that voids the
guarantee, and mitigated but not eliminated by backbone diversity; (iii)~compromise of
the trusted aggregation substrate or the instruction channel; and (iv)~attacks on the
underlying serving infrastructure. \Cref{fig:threat} maps each in-scope lever to the
component it targets and the defense that contains it.

% figures/tikz_threat_privacy_robustness.tex  ->  \cref{fig:threat}
\begin{figure}[t]
\centering
\resizebox{0.98\textwidth}{!}{%
\begin{tikzpicture}[node distance=9mm and 16mm]

% central pipeline spine (honest, blue) -------------------------------
\node[honest] (payload) {\icondoc[aegisBlue]\\ payload $\pay$\\ (data channel)};
\node[honest, right=of payload] (judges) {\iconjudge[aegisBlue]\\ honest judges\\ $\Hset$};
\node[honest, right=of judges] (agg) {\iconshield[aegisBlue]\\ aggregator\\ $\Agg$};
\node[honest, right=of agg] (dec) {\iconserver[aegisBlue]\\ decision $\hat d$};
\foreach \a/\b in {payload/judges, judges/agg, agg/dec}
  \draw[dataflow] (\a) -- (\b);

% attacker levers (red, top) ------------------------------------------
\node[adv, above=13mm of payload] (t2) {\iconattacker[aegisRed]\\ (ii) second-order\\ injection};
\node[byz, above=13mm of judges] (t1) {(i) compromise\\ up to $f$ judges};
\node[adv, above=13mm of agg] (t3) {(iii) adapt to\\ rule $\Agg$};

\draw[advflow] (t2.south) -- (payload.north) node[elabel,midway]{shift $\varepsilon$};
\draw[advflow] (t1.south) -- (judges.north) node[elabel,midway]{$f$ Byz.};
\draw[advflow] (t3.south) -- (agg.north) node[elabel,midway]{bias $\uhat$};

% defenses (green, bottom) --------------------------------------------
\node[defense, below=13mm of payload] (d2) {\iconlock[aegisGreen]\\ isolation\\ (\cref{prop:isolation})};
\node[defense, below=13mm of judges] (d1) {hardening\\ (SecAlign/StruQ)};
\node[defense, below=13mm of agg] (d3) {\iconshield[aegisGreen]\\ $2f{+}2<n$\\ (\cref{thm:integrity})};

\draw[defflow] (d2.north) -- (payload.south) node[elabel,midway]{contains};
\draw[defflow] (d1.north) -- (judges.south) node[elabel,midway]{resists};
\draw[defflow] (d3.north) -- (agg.south) node[elabel,midway]{bounds};

% legend --------------------------------------------------------------
\node[adv, minimum width=6mm, minimum height=4mm, below right=6mm and 6mm of dec] (lgA) {};
\node[right=1mm of lgA, aegissmall] (lgAt) {attack lever};
\node[defense, minimum width=6mm, minimum height=4mm, below=2mm of lgA] (lgD) {};
\node[right=1mm of lgD, aegissmall] {defense};
\node[honest, minimum width=6mm, minimum height=4mm, below=2mm of lgD] (lgH) {};
\node[right=1mm of lgH, aegissmall] {honest pipeline};

\end{tikzpicture}}
\caption{Threat-and-defense map. The honest pipeline spine (blue) runs left to right. Above
it (red), the three composable attacker levers: (i) compromising up to $f$ judges (dashed
Byzantine node), (ii) second-order injection of the payload, and (iii) adapting to the
aggregation rule. Below it (green), the mechanism that contains each lever: isolation for
(ii), judge hardening for the judges, and the $2f{+}2<n$ robust-aggregation guarantee for
(iii). Each attack and defense connects to the single component it acts on; lanes are
separated by colour and dash, and every label sits on a white chip.}
\label{fig:threat}
\end{figure}

